
This thesis is concerned with addressing two main questions: what is the expected lifetime of a dual spacecraft cooperative GNSS-R formation flying mission, and how accurate are its predictions adjusted for its underlying Basilisk framework prediction accuracy? To address these questions, two distinct simulators were developed, the Comparative Propagation Simulator, and the GNSS-R Formation Flight Simulator. In order to accurately estimate the formation lifetime, its operational constraints had to be modeled, including the satellite configuration, electronic power system, autonomous mode switching, attitude- and formation control. The estimated lifetime was predicted using modes derived from two simulation campaigns, studying the fuel consumption's sensitivity to the leader and follower deployment velocities, and the steady-state formation-maintenance fuel consumption, respectively. The final lifetime estimates derived from the combination of these two models were $413.6~\mathrm{days}$ and $366.5~\mathrm{days}$ for the most optimistic and pessimistic deployment cases, respectively. However, due to modelling uncertainty and observed formation drift, the estimate was evaluated to be conservative.


The Basilisk propagation configuration used in the lifetime analysis was developed as a continuation of the specialization project, where the previous exponential atmosphere model was replaced by the higher-fidelity NRLMSISE-00 atmosphere model. This thesis configuration was benchmarked against a continuously updating Skyfield-\gls{sgp4} baseline by evaluating both same-spacecraft state differences and same-formation drift differences. The same-spacecraft comparison showed that the thesis configuration diverged substantially less from the Skyfield-\gls{sgp4} reference than the specialization-project configuration. However, the same-formation comparison showed that the thesis configuration predicted larger passive leader-follower drift than the Skyfield-\gls{sgp4} baseline. Since the lifetime campaigns were therefore likely subject to higher relative drift than predicted by the comparison baseline, it supports the claim of conservative mission lifetime estimate.